\documentclass[12pt]{article}
\usepackage{graphicx} % Required for inserting images
\usepackage{enumitem}
\usepackage{parskip}
\usepackage[a4paper, margin=2.54cm]{geometry}
\usepackage{amssymb}
\usepackage{amsmath}
\usepackage{amsthm}
\newtheorem{definition}{Definition}[section]
\newtheorem{convention}{Convention}[section]
\newcommand{\fal}{\mathrel{\text{\reflectbox{$\Vdash$}}}}
\newcommand{\ver}{\Vdash}
\newcommand{\inever}{\mid \! \mid \!\!>}
\newcommand{\inefal}{< \! \! \mid \! \mid}
\newcommand{\K}{\mathcal{K}}
\newcommand{\SK}{\mathcal{S}^\mathcal{K}}
\newcommand{\ostar}{\circledast}
\usepackage{bm}
\usepackage{mathptmx}
\usepackage{setspace}
% --- APA CITATION SETUP ---
\usepackage{apacite}
\DeclareMathAlphabet{\mathcal}{OMS}{cmsy}{m}{n}
\bibliographystyle{apacite}
\singlespacing

\title{Truthmaker Semantical Approaches to Belief Revision}
\author{Shin JeongBean}
\date{\today}

\begin{document}
\setlength{\parskip}{8pt}
\setlength{\parindent}{1em} 
\maketitle
\begin{abstract}
    In this paper, I propose some frameworks for Belief Revision by using resources from Finean truthmaker semantics, resembling the traditional AGM approach. New frameworks will be able to decide a unique contraction, and typically lack some disputed features such as logical omniscience and recovery.
\end{abstract}
\section{Introduction}
In this paper, I will develop a family of hyperintensional frameworks for belief revision utilizing the resources provided by Kit Fine's truthmaker semantics. 
\par
The classical AGM framework (notably \citeauthor{Alchourron1985-ALCOTL-2}, 1985; the description below follows that of \citeauthor{vanDitmarsch2007-VANDEL-6}, 2008) represents an agent's doxastic state as a theory or belief set $\K$ in a propositional language. Let $\mathcal{L}_0$ be the propositional language generated by a countable set of atomic propositions $P$ under the classical truth-functional connectives. Let $\text{Cn}$ be the classical consequence operator, where $\text{Cn}(\Sigma) = \{ \phi \in \mathcal{L}_0 \mid \Sigma \vdash \phi \}$. A \emph{belief set} $\K \subseteq \mathcal{L}_0$ is defined as a set of formulas closed under $\text{Cn}$, i.e., $\K = \text{Cn}(\K)$. The unique inconsistent belief set is denoted by $\K_{\bot} = \mathcal{L}_0$.
\par
The AGM framework includes three basic operations that map a belief set $\K$ and an input sentence $\phi \in \mathcal{L}_0$ to a new belief set. Intuitively, 
\begin{enumerate}
    \item \textbf{Expansion ($\K \oplus \phi$):} The incorporation of new information $\phi$ into $\K$, without retracting any existing beliefs. 
    \item \textbf{Contraction ($\K \ominus \phi$):} The removal of a belief $\phi$ from $\K$, such that the resulting set no longer implies $\phi$, unless $\phi$ is a classical tautology ($\vdash \phi$).
    \item \textbf{Revision ($\K \ostar \phi$):} The integration of a new, potentially conflicting belief $\phi$ into $\K$ (if $\phi$ is consistent), while possibly some sentences in $\K$ should be removed so the resulted set is also consistent.
\end{enumerate}
\par
The specific behavior of expansion is characterized by the postulates $(\oplus 1) \text{ - } (\oplus 6)$ in Table \ref{tab:expansion}. This additionally captures the idea that $\K \oplus \phi = \text{Cn}(K \cup \{\phi\})$. 
\par
The postulates for expansion determine clear and unique result after expansion, whereas characterizing contraction is more complicate under the AGM approach. Here I will not go deeper beyond so called \textit{partial meet contraction}, which conforms to the postulates $(\ominus 1) \text{ - } (\ominus 6)$ in \ref{tab:pm contraction}. The idea behind it is roughly that if $\phi$ is contracted from $\K$, the resulted belief set is the intersection of the favored maximal subsets of $\K$ that does not imply $\phi$. If $\phi$ is not in $\K$, then the maximal subset of $\K$ that does not imply $\phi$ would be $\K$ itself. 
\begin{table}[htpb]
\centering
\begin{tabular}{lr}
\hline
$(\oplus 1)$ $\mathcal{K} \oplus \phi$ is a belief set & \textit{type} \\
$(\oplus 2)$ $\phi \in \K \oplus \phi$ & \textit{success} \\
$(\oplus 3)$ $\K \subseteq \K \oplus \phi$ & \textit{expansion} \\
$(\oplus 4)$ $\phi \in \K \implies \K = \K \oplus \phi$ & \textit{minimal action} \\
$(\oplus 5)$ For all $\mathcal{H}, \mathcal{K} \subseteq \mathcal{H} \implies \mathcal{K} \oplus \phi \subseteq \mathcal{H} \oplus \phi$ & \textit{monotony} \\
$(\oplus 6)$ $\K \oplus \phi$ is the smallest set satisfying $(\oplus 1) - (\oplus 5)$ & \textit{minimal change} \\
\hline
\end{tabular}
\caption{The AGM Postulates for Expansion.}
\label{tab:expansion}
\end{table}
\begin{table}[htpb]
\centering
\begin{tabular}{lr}
\hline
$(\ominus 1)$ $\mathcal{K} \ominus \phi$ is a belief set & \textit{type} \\
$(\ominus 2)$ $\mathcal{K} \ominus \phi \subseteq \mathcal{K}$ & \textit{contraction} \\
$(\ominus 3)$ $\phi \notin \mathcal{K} \implies \mathcal{K} = \mathcal{K} \ominus \phi$ & \textit{minimal action} \\
$(\ominus 4)$ $\not\vdash \phi \implies \phi \notin \mathcal{K} \ominus \phi$ & \textit{success} \\
$(\ominus 5)$ $\phi \in \mathcal{K} \implies \mathcal{K} \subseteq (\mathcal{K} \ominus \phi) \oplus \phi$ & \textit{recovery} \\
$(\ominus 6)$ $\vdash \phi \leftrightarrow \psi \implies \mathcal{K} \ominus \phi = \mathcal{K} \ominus \psi$ & \textit{extensionality} \\
\hline
\end{tabular}
\caption{The Basic AGM Postulates for Contraction.}
\label{tab:pm contraction}
\end{table}
\par
Once expansion and contraction is characterized, we can define revision in terms of these two operations by so-called \textit{Levi Identity}: 
\begin{enumerate}
    \centering
    \item[Levi Identity:] $K \ostar \phi = (K \ominus \neg\phi) \oplus \phi$
\end{enumerate}
The revision operation, defined with the expansion and the partial meet contraction conforming to the postulates above can be axiomatized via the postulates $(\ostar1) \text{ - } (\ostar6)$ in Table \ref{tab:revision}. 
\begin{table}[htpb]
\centering
\begin{tabular}{lr}
\hline
$(\ostar1)$ $\K \ostar \phi$ is a belief set & \textit{type} \\
$(\ostar2)$ $\phi \in \K \ostar \phi$ & \textit{success} \\
$(\ostar3)$ $\K \ostar \phi \subseteq \K \oplus \phi$ & \textit{upper bound} \\
$(\ostar4)$ $\lnot \phi \notin \K \implies \K \oplus \phi \subseteq \K \ostar \phi$ & \textit{lower bound} \\
$(\ostar5)$ $\K \ostar \phi = \K_{\bot} \iff \vdash \lnot \phi$ & \textit{triviality}\\
$(\ostar6)$ $\vdash \phi \leftrightarrow \psi \implies K \ostar \phi = K \ostar \psi$ & \textit{extensionality} \\
\hline
\end{tabular}
\caption{The AGM Postulates for Revision.}
\label{tab:revision}
\end{table}
\par
It would be redundant to spill more ink to praise the exquisity and the usefulness of the AGM approach, since it is already appreciated for several decades. Yet for some reason one might desire different kinds of framework. In section 1.1, I will propose the desiderata of the framework I will develop in this paper. Section 1.2 will be devoted to deploy the conceptual materials from Finean truthmaker semantics. Section 2 presents two simple frameworks that partially satisfies the desiderata in section 1.1. Several slightly more sophisticated frameworks which is stronger than those presented in \S 2 is developed in section 3, including which fully satisfies the desiderata with different degrees. Some general limitations of new frameworks are acknowledged in section 4, and it will be preceded by a formal appendix. 
\subsection{Desiderata}
\subsubsection{Definite Contraction}


The AGM framework assumes unique contraction. That is, by $(\ominus 1)$, the result of contracting $\phi$ from $\K$ must be one particular belief set. But the postulates $(\ominus 1) \text{ - } (\ominus 6)$ by themselves do not guarantee there is a unique belief set that satisfies them. 
\par
For example, suppose $\K = \text{Cn}(\{p, q, r\})$. In this case $\K \ominus (p \land q)$ has multiple possible outcomes. By $(\ominus 4)$ and the definition of belief set, it would be evident that at least one of $p$ and $q$ should be removed from the belief set; but which one? Either choice satisfies the postulates, and even if one removes all of them, it is still a valid partial meet contraction. But then, how do we choose one among them? In partial meet contraction, the choice was somewhat arbitrary. 
\par
The standard way out for the arbitrariness was by incorporating two more postulates $(\ominus 7)$ and $(\ominus 8)$ - which results in \textit{transitively relational partial meet contraction}. This turned out to capture the idea of \textit{epistemic entrenchment} \cite{gardenfors1988revisions}. In case where the epistemic agent is more entrenched into $p$, $\K \ominus (p \land q)$ would be , whereas in case it is more entrenched into $q$, the result would be Cn$({q, r})$. 
\begin{table}[htpb]
\centering
\begin{tabular}{lr}
\hline
$(\ominus 7)$ $((\K \ominus \phi) \cap (\K \ominus \psi)) \subseteq \K \ominus (\phi \land \psi)$ & \textit{min-conjunction} \\
$(\ominus 8)$ $\phi \notin \K \ominus (\phi \land \psi) \implies \K \ominus (\phi \land \psi) \subseteq \K \ominus \phi$ & \textit{max-conjunction} \\
\hline
\end{tabular}
\caption{Additional AGM Postulates for Contraction.}
\label{tab:entrenchment}
\end{table}
\par
Yet one could still complain that the entrenchment cannot be inferred from the belief set alone. The transitively relational partial meet contraction forces a regularity on contracting some formulas from $\K$, but still does not by itself decides what the regularity is. The regularity can be determined only if the entrenchment is already specified. But to fully specify the denumerably long list of entrenchment is, for one thing, practically not very realistic. Thus, one might desire a framework which can determine the result of contraction simply from what is contracted from what, just as expansion in the AGM framework is definite in that sense. 
\par
A related consideration is that one might want to model the epistemic entrenchment as a subject of change as the belief set undergoes some operations. Let $\K = \text{Cn}(\{p, \lnot q \})$. It seems natural to me to suppose that if $\K$ is revised with $p \land q$, not only $\lnot q$ is contracted, but also that $p$ is possibly in some sense \textit{reinforced}, so $p$ might be more entrenched to the agent after the operation of revision. Thus, one might desire entrenchment to be understood in terms of operations such as expansion and contraction; not to determine the result of contraction via entrenchment. 
\par
So this gives the first desideratum for frameworks I will develop in this paper:
\begin{enumerate}
    \item[(D1)] \textbf{Definite Contraction}: \textit{The definition, or the postulates of the contraction operator alone should be sufficient to determine the result after the contraction}
\end{enumerate}

\subsubsection{Logical Ignorance, Paraconsistency, and Local Rationality}
Logical omniscience is a traditional worry for any kinds of epistemic logic with the necessitation rule: $\vdash \phi \implies \vdash \Box \phi$. In epistemic logic, this gives that an epistemic agent have the doxastic or epistemic attitude which is supposed to conform to the behavior of $\Box$, towards all logical truths. Usually this yields the consequence that an ideal epistemic agent must know or believe all logical truths, which may seem suspicious. 
\par
The AGM approach yields the same result. Since the belief set $\K$ is defined as a set of formulas closed under the classical consequence operator, in any case the epistemic agents believes all the logical truths, which cause the same discomfort as the problem of logical omniscience for epistemic logics. 
\par
Additional worry is that the logical omniscience results in dismissing the contraction with logical truths. This is what \textit{minimal action} postulate commands; but even without the postulate the definition of $\K$ itself forbids the removal of any logical truths. Revising \textit{minimal action} with a weaker postulate at most only results in eccentric possibilities that by contracting a logical truth it rather takes out some innocent formulas other than what is contracted. However, one might want to model what an ideally rational epistemic agent must believe after it is forced to contract a logical truth from the belief set. 
\par
Likewise, one might feel uncomfortable to the idea that one must believe all the logical consequences of its belief in general. One case we might want to model is the pessimism for one's own belief. One might be equally entrenched deeply to each of the 'atomic' belief, but came to contract the conjunction of all of its beliefs, for it also has certainty towards at least some of its belief would be wrong; that the agent itself cannot be perfect. In this case, a faithful modelling seems not to remove all the beliefs (or all the logically intermediate beliefs) from the original belief set, but only to remove the conjunction, yet retaining all the `atomic' belief. But to this one must admit a belief set that is not closed under classical consequence. 
\par 
So this gives the second desideratum:
\begin{enumerate}
    \item[(D2)] \textbf{Logical Ignorance}: \textit{The belief set must be able to exclude its logical consequences}
\end{enumerate}
There is also a similar worry for expansion. Again, since $\K$ is closed under the classical consequence, starting with inconsistent beliefs, or expanding with any $\phi$ such that $\K \vdash \lnot \phi$ results in explosion. 
\par
But it may be unclear why an ideally rational epistemic agent should be so vulnerable to inconsistency. Especially, one might want to model how should one operate when the initial state of belief set is reasonably inconsistent. This gives the third desideratum:
\begin{enumerate}
    \item[(D3)] \textbf{Paraconsistency}: \textit{The belief set must be able to admit inconsistency without explosion}
\end{enumerate}
Finally, the motivating idea behind \textbf{Logical Ignorance} and \textbf{Paraconsistency} may be the consideration that we might want to model an epistemic agent that is \textit{locally} rational, in the sense that it contemplates a subset of the whole belief set at a time. Indeed, human epistemic agents never inspect every corner of our own belief set when expanding or contracting our beliefs; and this is somewhat desirable even for an ideal epistemic agent when one considers cognitive economy. 
\par
This is also motivated from the idea that the strength of those two desiderata must be somehow inhibited. For example, \textbf{Paraconsistency}, in its extreme, would allow one to believe apparent contradictions such as $\phi \land \lnot \phi$; yet we might want the modelled ideal epistemic agent be aware enough not to make such a silly mistake. \textbf{Logical Ignorance}, on the other hand, would make possible that the ideal agent believes both $\phi$ and $\lnot \phi \lor \psi \ (\iff \phi \rightarrow \psi)$ without believing $\psi$. Although we might be reluctant to take the modelled agent to be perfect in all logical regards, we still want it to be locally \textit{rational}. 
\par
So this gives the forth, somewhat vague desideratum:
\begin{enumerate}
    \item[(D4)] \textbf{Local Rationality}: \textit{The modelled epistemic agent must be rational, but only locally so.}
\end{enumerate}
\subsubsection{Hyperintensionality}
In \S 1.1.2, it was suggested that the new framework must be able to contract with logical truths. Naturally, if it is possible to contract with logical truths, it would seem also possible to expand, or revise with logical truths. 
\par
Yet with \textit{extensionality} postulates for contraction and revision, non-trivial contraction and revision with logical truths consist only of two kinds: either you contract \textit{all} logical truths, or you revise and thus incorporate \textit{all} logical truths. Suppose the belief set $\K$ such that  $\{p \lor \lnot p, q \lor \lnot q \} \subset \K$ and $r \lor \lnot r \notin \K$. This would be possible in frameworks where \textbf{Logical Ignorance} is satisfied. Now consider $\K \ominus (p \lor \lnot p)$. Since $\vdash (p \lor \lnot p) \leftrightarrow (q \lor \lnot q)$, by \textit{extensionality} of contraction $\K \ominus (p \lor \lnot p) = \K \ominus (q \lor \lnot q)$, and by \textit{success} we get $\{p \lor \lnot p, q \lor \lnot q \} \cap \K \ominus (p \lor \lnot p) = \varnothing$. Furthermore, by extensionality for revision, likewise we get $\{p \lor \lnot p, q \lor \lnot q, r \lor \lnot r \} \subset (\K \ominus (p \lor \lnot p)) \ostar (r \lor \lnot r)$. 
\par
But one might desire a more fine-grained contraction and revision. For one thing, if one already has \textbf{Local Rationality}, one might want to be able to contract some particular kind of logically true formulas without affecting other formulas, especially for those have different \textit{subject matter}. Likewise, one might want to be able to expand, or revise with some logically true formulas without bringing infinities of logical truths into the belief set. 
\par
So this gives the fifth desideratum:
\begin{enumerate}
    \item[(D5)] \textbf{Hyperintensionality}: two extensionally, or intensionally equivalent formulas must be able to behave differently upon expansion or contraction.
\end{enumerate}
\subsubsection{Pertinent Recovery}
Finally, it should be noted that the \textit{recovery} postulates was a subject of dispute for a long time. The idea that if one loses a belief and immediately regains it the agent must at least return to the original belief set would seem intuitively compelling. The postulate also plays a crucial role in preventing the contraction for being unnecessarily strong, so contracting with $\phi$ cannot remove more formulas from the belief set than what expanding with $\phi$ would take into the belief set. 
\par
Yet \textit{recovery}, as the AGM framework postulates, also yields counterintuitive results. Suppose $\{p, q, p \land q, p \lor q\} \subset \K$. Now, the fact that $\{p, q, p \land q, p \lor q\} \cap (\K \ominus (p\lor q) = \varnothing$ will follow in the AGM framework (by \textit{success} of contraction and the definition $\K = \text{Cn}(\K)$), and it is quite intuitive. If the modelled epistemic agent loses the belief that it is either the case that $p$ or $q$, reasonably the agent would also lose the belief that it is both the case that $p$ \textit{ and } $q$. So far so good, but \textit{recovery} yields that if the modelled agent immediately regains the belief that $p \lor q$, the agent must also regain the belief that $p \land q$. That is: $\{p, q, p \land q, p \lor q\} \subset (\K \ominus (p \lor q)) \oplus (p \lor q)$. But this is intuitively not very compelling.
\par
To generalize: suppose $\{\phi, \psi \} \subset \K$ where $\phi$ strictly entails $\psi$ but not the other way around. This yields the result that $\phi \notin \K \ominus \psi$. Yet by recovering $\psi$, the modelled agent recovers not only $\psi$ but the stronger $\phi$. This is also suspicious, for recovering $\psi$ seems to have no reason to involve $\phi$ if $\psi$ does not entail $\phi$, regardless of the fact that $\phi$ was in the initial belief set before the contraction \cite{sep-logic-belief-revision}.  
\par
It seems to me that the intuition behind proposed counterexamples for \textit{recovery} is that one might have different \textit{sources} for the same belief. If $\psi$ is to be contracted from $\K$, not only should the \textit{sources}, or reason for believing $\psi$ should be discarded, but also should the sources of the belief that $\phi$ be defeated, for any source of the belief that $\phi$ would also be the source of the belief that $\psi$. But when one recovers the belief that $\psi$, it is not guaranteed that the recovery is via reconfirming the exact source for both beliefs that $\psi$ and $\phi$. One may recover the belief that $\psi$ via \textit{new} source for the belief that $\psi$ which does not make one believe that $\phi$. It is possible because, since $\psi$ does not entail $\phi$, the source for the belief that $\psi$ is not necessarily the source for the belief that $\phi$. 
\par
Notice how \textit{recovery} is extremely compelling when it is formulated in terms of the source of belief: if a source of belief is once discarded and immediately reconfirmed, the original belief set must be a subset of the belief set after the sequence. Thus, one might wish the framework to be able to grasp this sense of recovery, not the one among the AGM postulates. 
\par
This gives the final, sixth desideratum:
\begin{enumerate}
    \item[(D6)] \textbf{Pertinent Recovery} The framework must avoid the problematic case of \textit{recovery}. 
\end{enumerate}
\subsection{Materials From Truthmaker Semantics}
In this paper, I propose frameworks that satisfy \textbf{Definite Contraction}, \textbf{Logical Ignorance}, \textbf{Paraconsistency}, \textbf{Hyperintensionality}, and \textbf{Source}. \textbf{Local Rationality} will not be satisfied in two prototypes, and will be realized in other frameworks in different degrees. 
\par 
The core idea is to posit a \textit{belief-state set} $\SK$, where any belief-state set determines a unique belief set. Intuitively, $\SK$ is the set of the \textit{source} of the modelled agent's belief; so any member of $\SK$ is a state which \textit{verifies}, or \textit{truthmakes} some member of $\K$. Operations are defined to operate on $\SK$ instead of $\K$; thus we expand, contract or revise with a state $\sigma$ from a belief-state set $\SK$. The expansion, contraction or revision with a formula $\phi$ from a belief set $\K$ is only defined in terms of the operation with state on a belief-state set. 
\par 
The ideas presented in this subsection originates from Kit Fine's work on truthmaker semantics. Below in this subsection I introduce some important resources from Finean truthmaker semantics which will be directly relevant to new frameworks developed in this paper. The following details are mostly taken from \citeA{Fine2016-FINAC-2}, \citeA{Fine2017-FINATO-3}, \citeA{Fine2017-FINATO-4}, and \citeA{Fine2017-FINTSP-2}.
\par
A state, informally, is anything that exactly verifies or exactly falsifies a statement (the notion of exact and inexact verification/falsification will be covered shortly). We do not need to assume a particular ontology of states to do semantics, as we do not need to assume a particular ontology of possible worlds to do modal logics. 
\par
One characteristic feature of states is that states can have mereological structure. Thus a state can be a part of another state; and any two states can fuse to compose a single state. The parthood relation here is supposed to be antisymmetric and transitive. 
\par
Formally, we achieve such an understanding on states by defining a \textit{set space}. To do so we first define the parthood relation $\sqsubseteq$ and the notion of \textit{least upper bound}.
\begin{definition}[Parthood Relation on $S$]
$\sqsubseteq$ (part) is a binary relation on a set of states $S$ which conforms to the following three conditions: \par
{\centering
Reflexivity: $\forall s (s \sqsubseteq s)$ \\
Antisymmetry: $\forall s \forall t ((s \sqsubseteq t \& t \sqsubseteq s) \implies s = t)$ \\
Transitivity: $\forall s \forall t \forall u ((s \sqsubseteq t \& t \sqsubseteq u) \implies s \sqsubseteq u)$ \\} 
\end{definition}
\begin{definition}[Least Upper Bound]
A state $s$ is an upper bound of $T \subseteq S$ iff $\forall t \! \in \! T \ t \sqsubseteq s$. \\
A state $s$ is a least upper bound of $T$ iff $s$ is an upper bound of $T$, and it is part of any upper bound $s'$ of $T$. 
\end{definition}
\begin{definition}[State Space]
A state space is an ordered set $\langle S, \sqsubseteq \rangle$, where $S$ is a set of states and $\sqsubseteq$(part) is a binary relation on $S$, and for any $T \subset S$ its least upper bound exists.
\end{definition}

It will be easy to check that the least upper bound is unique by antisymmetry. The least upper bound (sometimes called `fusion') of $T$ is denoted by $\bigsqcup T$, or $t_1 \sqcup t_2 \sqcup t_3 ...$ for $T = \{t_1, t_2, t_3, ...\}$). 
\par
Two extreme states are worth noting:
\begin{convention}[Null State and Full State]
Let `$\Box$' denote the least upper bound of $\varnothing \subseteq S$, which is a part of any state and has no part. This will also be called `null state'. Let `$\blacksquare$' denote the least upper bound of $S \subseteq S$, which is not a part of any state and has every state as part. This will also be called `full state'
\end{convention}

The definition for the propositional language we will use throughout the paper is defined as follows:
\par
\begin{definition}[Propositional Language $\mathcal{L}$]
    Let $P$ be the set of atomic formulas. The language $\mathcal{L}$ is generated by the following BNF: 
    \par
    {\centering $\phi$::= $\ \ p \  | \ \lnot \phi \ | \ (\phi \land \phi) \ | \ (\phi \lor \phi)$ \par}
    for any $p \in P$.
\end{definition}
We also deploy standard conventions for bracketing and definitions for other boolean connectives. Notably, $(\phi \rightarrow \psi)$ is defined as $(\lnot \phi \lor \psi)$.
\par
The \textit{truthamker model}, intuitively, decides a pair of sets - the set of verifiers and the set of falsifiers - of states in a state space for each atomic formulas. Given a model we can recursively define the verification and falsification condition for any formula in $\mathcal{L}$. Formally:
\begin{definition}[Modalized Set Space]
    A modalized set space is a triple $\langle S, S^\Diamond, \sqsubseteq \rangle$ where $\langle S, \sqsubseteq \rangle$ is a set space and $S^\Diamond$ (set of all possible states) is a non-empty subset of $S$ such that $(s \in S^\Diamond \ \& \ t \sqsubseteq s) \implies t \in S^\Diamond$.
\end{definition}
\begin{definition}[Model $\bm{M}$]
    A model $\bm{M}$ is a quadruple $\langle S, S^\Diamond, \sqsubseteq, | \cdot  | \rangle$ where $\langle S, S^\Diamond, \sqsubseteq \rangle$ is a modalized set space, and $| \cdot |$ is a valuation function which takes every $p \in P$ to a pair $\langle V, F \rangle$, where $V$ and $F$ are non-empty subset of $S$ that conforms to the following conditions: \par
    {\centering
    Exclusivity: $\forall v \! \in \! V \ \forall f \! \in \! F \ v \sqcup f \notin S^\Diamond$ \\
    Exhaustivity: $\forall s \! \in \! S^\Diamond (\exists v \! \in \! V \ s \sqcup v \in S^\Diamond \lor \exists f \! \in \! F \ s \sqcup f \in S^\Diamond)$ \par} 
    Given $|p| = \langle V, F \rangle $, let $|p|^+$ denote the first component (i.e. $V$), and $|p|^-$ denote the second component (i.e. $F$).
\end{definition}
Alternatively, we could allow $V$ or $F$ to be empty but never both. In fact, \textit{Exhaustivity} already implies that it can never be the case that both $V$ and $F$ are empty, and \textit{Exculsivity} grants that if one of them is empty, the other one would be $\{\Box\}$. 
\par
Given a model $\bm{M}$, we define a basic notion of verification and falsification, that is, \textit{exact} verification $\ver$ and \textit{exact} falsification $\fal$: 
\begin{definition}[Exact Verification and Falsification for Formuals]
Given a model $\bm{M} = \langle S, S^\Diamond, \sqsubseteq, | \cdot  | \rangle$, 
    \begin{enumerate}[align=left]
    \item[$(atom)^+$] $s \ver p \iff s \in |p|^+$ 
    \item[$(atom)^-$] $s \fal p \iff s \in |p|^-$ 
    \item[$(negation)^+$] $s \ver \lnot \phi \iff s \fal \phi$ 
    \item[$(negation)^-$] $s \fal \lnot \phi \iff s \ver \phi$
    \item[$(conjunction)^+$] $s \ver (\phi \land \psi) \iff \exists t, u \ s.t. \ t \ver \phi \ \& \ u \ver \psi \ \& \ s = t \sqcup u$
    \item[$(conjunction)^-$] $s \fal (\phi \land \psi) \iff s \fal \phi \text{ or } s \fal \psi$
    \item[$(disjunction)^+$] $s \ver (\phi \lor \psi) \iff s \ver \phi \text{ or } s \ver \psi$ 
    \item[$(disjunction)^-$] $s \fal (\phi \lor \psi) \iff \exists t, u \ s.t. \ t \fal \phi \ \& \ u \fal \psi \ \& \ s = t \sqcup u$
    \end{enumerate}
\end{definition}
Notice that the exact verifier of $\phi \land \psi$ does not always exactly verifies $\phi$. Likewise, the exact falsifier of $\phi \lor \psi$ does not always exactly falsifies $\phi$. Intuitively, a state $s$ exactly verifies $\phi$ only if $s$ is wholly relevant with the truth of $\phi$. For example, the state that it is raining is wholly relevant with the truth of the statement ``it is raining'', but the state that it is raining and there is a dog lying in the gutter is not wholly relevant with the truth of the statement ``it is raining''; it is merely \textit{partially} relevant. Similar with the case of exact falsification. 
\par
Yet we could say the state that it is raining and there is a dog lying in the gutter \textit{inexactly} verifies the statement ``it is raining''. With the mereological structure of states in hand, we can define \textit{inexact verification} $\inever$ and \textit{inexact falsification} $\inefal$ as follows: 
\begin{definition}[Inexact Verification and Falsification for Formuals] Given a model $\bm{M} = \langle S, S^\Diamond, \sqsubseteq, | \cdot  | \rangle$, 
    \begin{enumerate}[align=left]
    \item[$(inexact)^+$] $s \inever \phi \iff \exists t \sqsubseteq s \ s. t. \ t \ver \phi$ 
    \item[$(inexact)^-$] $s \inefal \phi \iff \exists t \sqsubseteq s \ s. t. \ t \fal \phi$ 
    \end{enumerate}
\end{definition}
Finally, we define the notion of \textit{subject matter}: 
\begin{definition}[Subject Matter] 
    Given a model $\bm{M} = \langle S, S^\Diamond, \sqsubseteq, | \cdot  | \rangle$, let $\mathcal{V}(\phi)$ denote the set of all exact verifiers of $\phi$; and $\mathcal{F}(\phi)$ denote the set of all exact falsifiers of $\phi$. \par
    The subject matter of $\phi$ $sm(\phi)$ is a state in $S$ such that $\bigsqcup \mathcal{V}(\phi) \sqcup \bigsqcup \mathcal{F}(\phi)$.\footnote{
        The standard definition of subject matter is in terms of \textit{propositions} defined in terms of states: a \textit{bilateral proposition} $P$ is defined as $\langle P^+, P^- \rangle$ where $P^+$ and $P^-$ contains its verifiers and falsifiers respectively. Then the subject matter of $P$ is either the pair $\langle \bigsqcup P^+, \bigsqcup P^- \rangle$ (\textit{differentiated subject matter}) or the fusion of $\bigsqcup P^+$ and $\bigsqcup P^-$ (\textit{comprehensive subject matter}). I intentionally deviated from the standard definition for the project of this paper the commitment for the theory of turthmaker content is not required. The preference for the comprehensive over the differentiated subject matter in epistemic settings follows \citeA{Jago2024-JAGBIT}.
    }
\end{definition}


\section{Simple Frameworks}
With the materials from truthmakers semantics, now I develop two simple frameworks $\mathcal{E}_0$ and $\mathcal{I}_0$ that immediately arise from the idea that the expansion and contraction operates on a belief-state set $\SK$. Two simple frameworks will only differ in how to determine the belief set $\K$ from $\SK$. Thus the definitions in this section apply to both frameworks if not otherwise noted. 
\par
In simple frameworks we do not pose any constraint on $\SK$:
\par
\begin{definition}[Simple Belief-State Set]
    Given a model $\bm{M} = \langle S, S^\Diamond, \sqsubseteq, | \cdot  | \rangle$, a belief-state set $\SK$ is a subset of $S$. \par
\end{definition}
\par
In simple frameworks, an expansion of $\SK$ with a state $\sigma$ ($\SK + \sigma$) simply adds $\sigma$ to $\SK$ if not already contained; and a contraction of $\SK$ with $\sigma$ ($\SK \div \sigma$) simply extracts $\sigma$ from $\SK$ only if it was originally contained. Thus:
\begin{definition}[Expansion and Contraction in Simple Frameworks]
    Given a belief-state set $\SK$, $\SK + \sigma = \SK \cup \{\sigma\}$; and $\SK \div \sigma = \SK \setminus \{\sigma \}$.
\end{definition}


\par
As the ontology of states does not matter for our purpose, the ontology of belief-state set is also not required. Yet intuitively, one can regard $\SK$ as a set of states that the modelled epistemic agent is epistemically committed into; that is, the set of sources of all its beliefs. My preferred understanding is to take $\SK$ as a model for the modelled epistemic agent's `representation' of reality, in the broadest sense. Thus, the belief set would contain every formula that would be actually verified if the representation were totally correct. Since we have two different notions of verification: exact and inexact, we would have two different frameworks that decide $\K$ from $\SK$:
\begin{definition}[Belief Set in $\mathcal{E}_0$ and $\mathcal{I}_0$] In framework $\mathcal{E}_0$, given a simple belief-state set $\SK$, a belief set $\K = \{\phi : \exists \sigma\in\SK \ s.t. \ \sigma \ver \phi \} \subseteq \mathcal{L}$. 
\par
In framework $\mathcal{I}_0$, given a simple belief-state set $\SK$, a belief set $\K = \{\phi : \exists \sigma\in\SK \ s.t. \ \sigma \inever \phi \} \subseteq \mathcal{L}$. 
\end{definition}
Notice that two basic frameworks only differ in this respect. We could define a framework as a $n$-tuple. Thus:
\begin{definition}[Frameworks $\mathcal{E}_0$ and $\mathcal{I}_0$]
    Given a model $\bm{M}$, the framework $\mathcal{E}_0$ is a quintuple $\langle \bm{M}, \Sigma_0, +, \div, \Delta_{\mathcal{E}} \rangle$ where $\Sigma_0$ is a set of belief-state set $\{\SK: \SK \subseteq S\}$, $+$ and $\div$ are expansion and contraction respectively, and $\Delta$s are decision function  from the set of belief-state sets to the set of belief sets such that $\Delta_{\mathcal{E}} (X) = \{\phi : \exists \sigma\in X \ s.t. \ \sigma \ver \phi \} \subseteq \mathcal{L}$.
    \par
    Given a model $\bm{M}$, the framework $\mathcal{I}_0$ is a quintuple $\langle \bm{M}, \Sigma_0, +, \div, \Delta_{\mathcal{I}} \rangle$ where $\Delta_{\mathcal{I}} (X) = \{\phi : \exists \sigma\in X \ s.t. \ \sigma \inever \phi \} \subseteq \mathcal{L}$, and everything else is as defined above. 
\end{definition}
\par
One might desire there to be additional notions of expansion and contraction that operates on $\K$. Intuitively, the expansion of $\K$ with a formula $\phi$ $(\K \oplus \phi)$ would be the resulting belief set of the belief-state set $\SK$ expanded with \textit{some} verifiers of $\phi$; the contraction of $\K$ with a formula $\phi$ ($\K \ominus \phi$) would be the resulting belief set of the belief-state set $\SK$ contracted with \textit{all} verifiers of $\phi$. 
\par
However, it should be noticed that given a definite $\SK$, the resulting belief set $\K$ is definite in both $\mathcal{E}_0$ and $\mathcal{I}_0$; whereas given a definite belief set one cannot track down to a unique belief-state set. So let us denote the particular belief set decided by a belief-state set $\SK$ $D(\SK)$. More specifically, the belief state decided by a belief-state set $\SK$ \textit{in a framework X} will be denoted by $D_X(\SK)$. The specification of the framework will be omitted unless context is insufficient to disambiguate. Then, the expansion and contraction for belief set should be defined as follows:
\par
\begin{definition}[Operations on Belief Set in Simple Frameworks]
    Given a belief-state set $\SK$ and a selection function $\gamma$,
    \begin{enumerate}
        \item[($\mathcal{E}_0^\oplus$)] $D_{\mathcal{E}_0} (\SK) \oplus \phi = \gamma (D_{\mathcal{E}_0} (\SK + \sigma) \ s.t. \ \sigma \ver \phi$ 
        \item[($\mathcal{E}_0^\ominus$)] $D_{\mathcal{E}_0} (\SK) \ominus \phi = D_{\mathcal{E}_0} (\SK \setminus \{\sigma: \sigma \ver \phi \})$ 
        \item[($\mathcal{I}_0^\oplus$)] $D_{\mathcal{I}_0} (\SK) \oplus \phi = \gamma(D_{\mathcal{E}_0} (\SK + \sigma)) \ s.t. \ \sigma \ver \phi$ 
        \item[($\mathcal{I}_0^\ominus$)] $D_{\mathcal{I}_0} (\SK) \ominus \phi = D_{\mathcal{E}_0} (\SK \setminus \{\sigma: \sigma \inever \phi \})$ 
    \end{enumerate}
    The revision on belief set ($\ostar$) is defined via Levi Identity: 
    \begin{enumerate}
        \item[($\mathcal{E}_0^\ostar$)] $D_{\mathcal{E}_0} (\SK) \ostar \phi = (D_{\mathcal{E}_0} (\SK) \ominus \lnot \phi) \oplus \phi$
        \item[($\mathcal{I}_0^\ostar$)] $D_{\mathcal{I}_0} (\SK) \ostar \phi = (D_{\mathcal{I}_0} (\SK) \ominus \lnot \phi) \oplus \phi$
    \end{enumerate}
\end{definition}
This yields the result that still expansion, and thus revision on belief set is not definite in $\mathcal{E}_0$ nor in $\mathcal{I}_0$, for a formula can have multiple verifiers, exact or inexact. But I got an excuse for that. As I said, we can think $\SK$ to be the set of all doxastic sources, or reason the modelled epistemic agent have upon forming beliefs. So it seems natural to me that, given a belief-state set, one cannot determine what would be the result after one acquired a belief, for it is not specified what is the exact source of that acquired belief. However, the definition also allows the expansion infringe \textit{minimal action} for expansion (see Table \ref{tab:expansion}). This feature, I admit, may be disputed. One might add the clause ``if $\phi \in D(\SK)$, $D(\SK) \oplus \phi = D(\SK)$'' to the definition to retain the minimal action postulate, if one whishes. On the contrary, it seems given a belief-state set, we could determine what would be the result after one discarded a belief: the agent would have lost \textit{all} the sources for the discarded belief. Also notice that in case of contraction, \textit{minimal action} for contraction (see Table \ref{tab:pm contraction}) holds. 
\par
Now let us observe how two simple frameworks satisfies every desideratum except \textbf{Local Rationality}. 
\begin{enumerate}
    \item[(D1)] \textbf{Definite Contraction}
\end{enumerate}
It would be easy to check that the contraction (on belief-state set) is uniquely determined, since relative complement is uniquely determined. 
\par
One interesting feature of these frameworks is that we can now define \textit{entrenchment} of the belief that $\phi \in D(\SK)$ by inspecting $\SK$. As for verification, we have exact and inexact entrenchment, such that:
\begin{definition}[Entrenchment]
    The exact entrenchment of $\phi \in D(\SK)$ is the number of exact verifiers of $\phi$ in $\SK$, i.e., $\text{card}(\SK \cap \{\sigma: \sigma \ver \phi\})$ \par
    The inexact entrenchment of $\phi \in D(\SK)$ is the number of inexact verifiers of $\phi$ in $\SK$, i.e., $\text{card}(\SK \cap \{\sigma: \sigma \inever \phi \})$
\end{definition}
The intuitive idea behind is that the modelled epistemic agent is more entrenched to the belief that $\phi$ just in case it has more sources or reason of believing that $\phi$. Observe that different frameworks get along well with different notions of entrenchment: the natural notion of entrenchment in framework $\mathcal{E}_0$ is \textit{exact entrenchment}, while entrenchment in framework $\mathcal{I}_0$ it is \textit{inexact entrenchment}.
\par
With the notion of entrenchment, one thing we can now make sense is the phenomenon that sometimes expanding or contracting on $\SK$ which makes difference on the belief-state set may still not result in any difference on $\K$. Suppose the state $t \notin \SK$ only verifies $\phi$, but $s \in \SK$ such that $s \ver \phi$, and all fusions of $t$ and any state in $\SK$ is not in $\SK$. Then $\SK + t$ would make no difference in $\K$, under both frameworks $\mathcal{E}_0$ and $\mathcal{I}_0$; likewise $\SK \cup \{ t \} \div t$ would make no difference in $\K$. This, I submit, can be made sense by entrenchment. Although the belief set remains the same, by expanding with $t$ the entrenchment of $\phi$ has increased, i.e., the belief that $\phi$ is reinforced. Likewise, by contracting with $t$ the entrenchment of $\phi$ has decreased, i.e., the belief that $\phi$ is weakened. 
\begin{enumerate}
    \item[(D2)] \textbf{Logical Ignorance}
\end{enumerate}
This is relatively easy to show. In fact, it is rather difficult to construct a belief-state set that makes all logical truths contained in the belief set. To establish logical omniscience, the belief-state set must include, for any formula, either the verifier of it or the falsifier of it. Otherwise there are some logical truths that is not a member of $\K$. One extreme example is where $\SK = \varnothing$, which is allowed under the definition. Then by definition for belief set, $\K$ would contain no formula at all, needless to say all logically true formulas. 
\begin{enumerate}
    \item[(D3)] \textbf{Paraconsistency}
\end{enumerate}
The fact that the frameworks satisfy \textbf{Paraconsistency} is also very much direct. Notice that in simple frameworks there are no constraint on $\SK$ that it should not contain both verifiers and falsifiers for some formulas. Thus, suppose $\SK = \{s, t\}$ where $s$ verifies (exactly for $\mathcal{E}_0$, inexactly for $\mathcal{I}_0$. Similarly, for below) $\phi$ and $t$ falsifies $\phi$ (therefore verifies $\lnot \phi$); and each state in $\SK$ verifies or falsifies nothing else. Then $D(\SK) = \{\phi, \lnot \phi \}$, without introducing all formulas in $\L$ into $D(\SK)$. 
\par
However, this maybe thought also as a weakness of simple frameworks, for it means that the frameworks allow blatant contradictions. The details will be discussed below when discussing \textbf{Local Rationality}.
\begin{enumerate}
    \item[(D5)] \textbf{Hyperintensionality}
\end{enumerate}
Hyperintensionality is automatically granted in our frameworks, for underlying truthmaker model is already hyperintensional. Suppose $s_1 \ver p$ and $s_2 \fal p$; $t_1 \ver q$ and $t_2 \fal q$. What follows is that although $\vdash p \lor \lnot p$ and $\vdash q \lor \lnot q$, they may have different verifiers: for $s_1 \ver p \lor \lnot p$ and $t_1 \ver q \lor \lnot q$. Likewise, although $p \land \lnot p$ and $q \land \lnot q$ are all logical falsities, they may have different verifiers: for $s_1 \sqcup s_2 \ver p \land \lnot p$ and $t_1 \sqcup t_2 \ver q \land \lnot q$. Recall that there were no constraints on model $\bm{M}$ that any states must verify logical truths or falsify logical falsities. Thus expanding or contracting with $p \lor \lnot p$ on $\K$ is not equivalent with expanding or contracting with $q \lor \lnot q$; likewise expanding or contracting with $p \land \lnot p$ on $\K$ is not equivalent with expanding or contracting with $q \land \lnot q$. 
\begin{enumerate}
    \item[(D6)] \textbf{Pertinent Recovery}
\end{enumerate}
Recall that the problematic case of recovery: $\phi$ strictly implies $\psi$, so contracting with $\psi$ would result in removing $\phi$ from $\K$; yet by \textit{recovery} (in Table \ref{tab:pm contraction}) expanding with $\psi$ would recover not only $\psi$ but also $\phi$. However, this contradicts with our daily observations, for the modelled epistemic agent could believe $\psi$ with some sources that do not force it to believe $\phi$. 
\par
The proposed frameworks deal with this kind of problem particularly well. Let us suppose all verifiers of $\phi$, including $s$, are also verifiers of $\psi$, but there is a state $t$ which verifies $\psi$ but not $\phi$. By supposition, if $s \in \SK$, $\K$ will include both $\phi$ and $\psi$; and $\{\phi, \psi\} \cap (D(\SK) \ominus \psi) = \varnothing$. But still, it is not necessarily that $\phi \in (D(\SK) \ominus \psi) \oplus \psi$. Since expansion on a belief set is not uniquely determined, $(D(\SK) \ominus \psi) \oplus \psi$ might be the set that contains the state $s$ that verifies both $\phi$ and $\psi$, but might as well be the set that contains the state $t$ that does not verify $\phi$, and thus $\phi \notin (D(\SK) \ominus \psi) \oplus \psi$. 
\par
So the idea that, given merely the assumptions that $\phi$ entails $\psi$ and $\{\psi, \phi\} \subseteq \K$ we cannot know whether $(\K \ominus \psi) \oplus \psi$ includes $\phi$ is preserved. We of course, can know whether $s$ is in $(\SK \div t) + t$ given the information whether $\{s\} \subseteq \SK$. But as noted in \S 1.1.4, the recovery formulated in terms of the source of belief - states and belief-state set - is extremely compelling. 
\par
To sum up, two simple frameworks - $\mathcal{E}_0$ and $\mathcal{I}_0$ already do very well in satisfying all the desiderata except \textbf{Local Rationality}. Yet they are untenable for they fail to satisfy the only remaining desideratum so drastically. 
\subsection{Rationality} 
As mentioned above, because there were no constraints on $\SK$ in both simple frameworks, they allow $\{\phi, \lnot \phi\} \subseteq \K$, and also $\phi \land \lnot \phi \in \K$ without any tension. Although I do not submit that the modelled agent should never be capable to have inconsistent belief set, I do hold that inconsistent belief set must cause some sort of `tension'. This is required if one wants to have a suitable notion of revision on $\SK$. 
\par
Recall that in simple frameworks revision is defined only on a belief set. This is because without any constraints on consistency of $\SK$, any attempts to define revision would be arbitrary, or too strong to satisfy other desiderata. For instance, one might wish to define revising with $\sigma$ on $\SK$ $(\SK * \sigma)$ to be $(\SK \setminus \{\tau: \tau \sqcup \sigma \notin S^\Diamond\}) \cup \{\sigma\}$. The idea is that, if we were to revise with $\sigma$, we first contract with all states that is \textit{incompatible} with $\sigma$, and then expand with $\sigma$. Then by definition of a model, $D(\SK * \sigma)$ will exclude all the formulas that is exactly falsified by $\sigma$. So far so good, but this also excludes all formulas in $D(\SK)$ such that every verifier of it is an impossible state; i.e., a logical falsity, for $\tau \sqcup \sigma \notin S^\Diamond$ if $\tau \notin S^\Diamond$. So after a revision on belief-state set with $\sigma$ it would result in removing all logical falsities. This would infringe our desideratum on \textbf{Hyperintensionality}. What if $\sigma$ already is an impossible state? Actually the result is much more devastating: if $\sigma \notin S^\Diamond$, according to the proposed definition on revision, $\SK * \sigma = \{\sigma\}$, for any $\tau$ will fuse with $\sigma$ to an impossible state. Since we wanted the modelled agent to have minimal changes even it is forced to believe inconsistently, we might want to say this violates \textbf{Paraconsistency}. 
\par
Another unfortunate feature of simple frameworks is that \textit{every} belief must be incorporated to the agent manually. $\mathcal{I}_0$ is perhaps better in this respect for by believing a conjunction it forces the agent to believe each of its conjuncts, while $\mathcal{E}_0$ does not yield this feature since the exact verifier of a conjunction is not generally the exact verifier of its conjuncts. But it should be examined whether this feature is really desirable. I postpone this discussion to \S 2.3.
\par
One especially problematic consequence of both basic frameworks is that one can believe that $\phi$ and $\phi \rightarrow \psi (\iff \lnot \phi \lor \psi$, since $\rightarrow$ is not a primitive operator in our language), yet not believing $\psi$. For the reason I will suggest in the next subsection, I admit there can be cases where one believes both $\phi \land \chi$ and $\phi \rightarrow \psi$, but not $\psi$. But it seems to me, a locally rational epistemic agent must believe that $\psi$, given that the agent believes that $\phi$, and it also believes that either it is the case that $\lnot \phi$ or $\psi$. In short, we want the modelled epistemic agent rational enough to have a belief set closed under modus ponens. 
\par
To sum up, simple frameworks fail to satisfy the `rationality' part of \textbf{Local Rationality}, for it permits blatant inconsistency in $\K$, and also because it fails to make the belief set closed under modus ponens (of material conditional). One might additionally complain that it is not closed under conjunction elimination. It is, for some parts, too weak to satisfy the desideratum. However, it is sometimes too strong to satisfy the `locality' part of \textbf{Local Rationality}. 
\subsection{Locality}
In some sense, it may be said that to model the locality of the rationality of the modelled epistemic agent is more important to our present purpose. For, if one only wanted the belief set to be gently consistent, or closed under modus ponens or conjunction elimination, the traditional AGM is already sufficient. In the previous section, we observed some cases that the frameworks fail to yield even local rationality. But there are also cases where the frameworks yield consequences than what we might desire.
\par
One crucial failure of locality is that in both simple frameworks, the belief set is closed under disjunction introduction: if $\phi \in \K$, for any formula $\psi$, $\phi \lor \psi \in \K$. It is easy to notice that we do not believe this way: although Frege would have believed that Wilhelm II is the King of Prussia, he would have not believed that Wilhelm II is the King of Prussia or Friedrich Merz is the chancellor of Germany. 
\par
Besides of intuitive implausibility, the closure under disjunction introduction also has further problems. For one thing, this makes any modelled epistemic agent believe \textit{about} everything. Let us define the subject matter of a \textit{set} of formulas $X$ as the least upper bound of the set of all subject matter of the formulas in $X$, formally:
\begin{definition}[Subject Matter for Sets]
    Given a model $\bm{M} = \langle S, S^\Diamond, \sqsubseteq, | \cdot  | \rangle$, the subject matter of a set $X \subseteq \mathcal{L}$ is defined as $sm(X) = \bigsqcup \{\sigma: \exists \phi \! \in \! X \ s. t. \ sm(\phi) = \sigma\}$.
\end{definition}
With the notion of subject matter for sets, we might want to investigate what the modelled epistemic agent believes \textit{about}. One intriguing example is that we can now define a stronger notion of ignorance: the agent is ignorant about $\phi$ if $\{\phi, \lnot \phi\} \cap \K = \varnothing$; but the agent is strongly ignorant about $\phi$ if the only common part of $sm(\phi)$ and $sm(\K)$ is $\Box$. 
\par
Yet if $\K$ is closed under disjunction introduction, there are no formulas that the agent is strongly ignorant. In fact, for any possible $\K$ that is not empty, $sm(\K) = sm(\mathcal{L})$ under both simple frameworks. This is clearly not desirable. It seems that until I searched on the internet to write this paragraph, I was strongly ignorant about that Amr ibn al-Layth was an Amir of Saffarid dynasty. But simple frameworks deny this intuition, for since I believe that I have a hand, my belief-state set would include a state $s$ that verifiers the statement ``I have a hand'', which also verifies ``I have a hand or Amr ibn al-Layth was an Amir of Saffarid dynasty''. And since the subject matter of the statement ``Amr ibn al-Layth was an Amir of Saffarid dynasty'' is a part of the subject matter of the disjunction, I would be not strongly ignorant about the statement about Amr ibn al-Layth, which I never have heard before. 
\par
Thus, it would be better to have a belief set that is not closed under disjunction introduction. However, excluding the disjunctive consequence of beliefs should be carefully executed; for we do not wish to ban the modelled agent to have disjunction as its belief. Notice that any verifiers, exact or inexact, of a disjunction $\phi \lor \psi$ must be a verifier, exact or inexact, of at least one of its disjuncts $\phi$ or $\psi$. So if we restrict the modelled epistemic agent believe $\phi \lor \psi$ on the basis of the belief that $\phi$, the result is that none of the member of $\K$ is a disjunction. This suggests that we may need a more fine-grained constraint on locality.


\subsection{Closure Under Conjunction Elimination?}
Finally, let me also note a problem about conjunction. Notice that in $\mathcal{E}_0$ even if $\phi \land \psi \in \K$, it does not grant that $\{\phi, \psi\} \subseteq \K$, whereas in $\mathcal{I}_0$ it is guaranteed. So at the first sight, it may seem as a reason to favor $\mathcal{I}_0$ over $\mathcal{E}_0$. One might therefore conclude that $\K$ should be defined from $\SK$ via the notion of \textit{inexact} verification rather than the notion of \textit{exact} verification. 
\par
The idea that the belief set should be closed under conjunction elimination seems plausible. However, I do not wish to adjudicate which notion of verification is suitable for defining $\K$ from $\SK$. The worry of using the exact notion of verification is as stated: the belief set can contain a conjunction without containing any of its conjuncts, which might make the modelled agent seem less rational than we wish it to be. Indeed, there might be something wrong with an agent if it believes $\phi \land \psi$ without believing that $\phi$ or $\psi$. 
\par
But using the inexact notion of verification is not without problems. We say $\phi$ is a conjunctive part of $\psi$ (and $\psi$ is a conjunctive whole of $\phi$) iff: (a) for all $s \ver \psi$, there is a $t$ such that $t \sqsubseteq s$ and $t \ver \phi$, and (b) for all $t \ver \phi$, there is a $s$ such that $t \sqsubseteq s$ and $s \ver \psi$ (and (c) any falsifier of $\phi$ is a falsifier of $\psi$. Clause (c) may be omitted sometimes)\footnote{Again, the notion of conjunctive part is also usually defined in terms of propositions; but let us operate with a watered-down notion of conjunctive part of formulas.}. Notice that the definition for conjunctive part does not involve the syntactic structure of $\phi$ or $\psi$. Now, under $\mathcal{I}_0$ not only each conjunct of a conjunction in $\K$ is included in $\K$, but all \textit{conjunctive part} of a member of $\K$ is included in $\K$. This is direct from the fact that any inexact verifier of a conjunctive whole would be an inexact verifier of its conjunctive part. 
\par
But it is unclear why should the modelled epistemic agent must believe all its conjunctive parts of its belief, especially when the conjunctive part in question is not syntactically related to its whole. At least in some cases, it seems there are good reason to believe that believing the conjunctive whole give sufficient condition to believe its conjunctive part regardless of their syntax. For example, it is reasonable to think ``$x$ is red'' is a conjunctive part of ``$x$ is scarlet''; and if one believes that $x$ is scarlet, one has sufficient reason to believe that $x$ is red \cite{Jago2024-JAGBIT}. 
\par
But it should also be noticed that sometimes the notion of conjunctive parthood are used to formulate semantic or metaphysical notions. It is customary to think if $\phi$ is a conjunctive part of $\psi$, the proposition that $\phi$ is a \textit{partial content} of the proposition that $\psi$ (Fine 2016, 2017a, 2017b). And it is also suggested that $\phi$ is a strict partial ground of $\psi$ iff $\phi$ is a conjunctive part of $\psi$ but not vice versa \cite{Fine2012-FINTPL-2}; or that $\phi$ a partial essence of $\psi$ iff $\phi$ is a conjunctive part of $\psi$ \cite{Correia2017-CORGEA}. If these suggestions are correct, defining $\K$ to contain every formula inexactly verified by some members of $\SK$ would result in the modelled agent to believe every partial contents, every strict partial grounds, and every partial essence of its beliefs - which sounds very suspicious. Of course, none of these suggestions are unchallenged; but $\mathcal{I}_0$ or its kin would be just as suspicious as those accounts are compelling. 
\par
Still, I personally favor to define $\K$ from $\SK$ by inexact verification. At least the semantic application of conjunctive parthood does not seem to threaten $\mathcal{I}_0$, for it may be claimed reasonable to expect an ideal epistemic agent should not believe $\phi$ if it is not fully aware of its semantic content. Nevertheless, the consideration above shows that frameworks that define $\K$ from $\SK$ by exact verification also deserve substantial interest. Also, as we will see in the next section, the approach also validates more fine-grained distinction on the constraints on $\K$. For such reasons both approaches will be developed in the remainder of the paper.
\par
A final remark. There is in fact a moderatist alternative between exact and inexact approaches: given a model $\bm{M}$ and a belief-state set $\SK$, we first define $\K_0$ to be the set of all formulas \textit{exactly} verified by some members of $\SK$; $\K_{n + 1}$ is defined as $\K_n \cup \{\phi, \psi: \phi \land \psi \in K_n\}$; and finally the belief set $\K$ is defined as $\bigcup_{n=0}^{\infty} \K_n$. This yields the result that $\K$ is closed under conjunction elimination, but not necessarily closed under conjunctive parthood. I will not develop this kind of approach further in this paper, but this is indeed an option. 

\section{More Frameworks!}
\subsection{The Locality Constraint}
Let us solve the problems of simple frameworks step-by-step. It turns out that the problem of disjunction is much easier to be dealt with. We only modify the definition of belief set to get the simple frameworks supplemented with the locality constraint:
\begin{definition}[Framework $\mathcal{ED}_0$]
    Given a model $\bm{M}$, the framework $\mathcal{ED}_0$ is a quintuple $\langle \bm{M}, \Sigma_0, +, \div, \Delta_{\mathcal{ED}} \rangle$ where $\Delta_{\mathcal{ED}} (X) = \{\phi : \exists \sigma\in X \ s.t. \ \sigma \ver \phi \ \& \ sm(\phi) \sqsubseteq \bigsqcup X\} \subseteq \mathcal{L}$, and everything else is as defined in \emph{\textbf{Definition 2.4}}. 
\end{definition}
\begin{definition}[Framework $\mathcal{ID}_0$]
    Given a model $\bm{M}$, the framework $\mathcal{ID}_0$ is a quintuple $\langle \bm{M}, \Sigma_0, +, \div, \Delta_{\mathcal{ID}} \rangle$ where $\Delta_{\mathcal{ID}} (X) = \{\phi : \exists \sigma \in X \ s.t. \ \sigma \inever \phi \ \& \ sm(\phi) \sqsubseteq \bigsqcup X\} \subseteq \mathcal{L}$, and everything else is as defined in \emph{\textbf{Definition 2.4}}. 
\end{definition}
The expansion and contraction with formulas under the framework $\mathcal{ED}_0$ is defined as it was defined under $\mathcal{E}_0$; and similarly under $\mathcal{ID}_0$ with $\mathcal{I}_0$. 
\par
The only differences of these two new frameworks are that we have an additional locality constraints on deciding $\K$ from $\SK$: intuitively, the subject matter of formulas in $\K$ must not exceed the ``scope'' of $\SK$. Given $\phi \in \K$, the locality constraint yields that $\phi \lor \psi \in \K \iff sm(\psi) \sqsubseteq \bigsqcup \SK$; thus if $\SK$ does not even contain a part of a state that exactly verifies or falsifies $\psi$ except $\Box$, $\phi \lor \psi \notin \K$ although two frameworks do not forbid $\K$ to include some disjunctions. Roughly, by incorporating the locality constraint on $\Delta$, we allow the modelled epistemic agent to have disjunctive belief, but only to the extent that each of the disjuncts are ``partially entertained in the representation'', positively or negatively. 
\par
This restores the intuition that epistemic agents can be strongly ignorant to some formulas. Suppose $\SK = \{s_1, s_2\}$, $|p|^+ = \{t_1\}$ and $|p|^- = \{t_2\}$, and $s_1 \ver q$ and $s_2 \fal r$. Finally, suppose $sm(p) (= t_1 \sqcup t_2) \not\sqsubseteq \bigsqcup \SK (= s_1 \sqcup s_2)$. Then, under both frameworks $\mathcal{ED}_0$ and $\mathcal{ID}_0$, $\{q \lor r, q \lor \lnot r\} \subseteq \K$, but $\{q \lor p, q \lor \lnot p\} \cap \K = \varnothing$. Furthermore, if $q$ nor $r$ has its exact verifiers or falsifiers which has $t_1$ or $t_2$ as part, and there are no other atomic formulas, we could conclude that the agent which has such a belief set is strongly ignorant about $p$, i.e., the only common part of $sm(p)$ and $sm(\K)$ is $\Box$. 
\par
Notice that in both frameworks, every discussed feature of simple frameworks remain: it decides a definite contraction, it satisfies \textbf{Logical Ignorance} - perhaps even better, for now the agent can be ignorant that from its belief $\phi$ it follows that $\phi \lor \psi$ for some $\psi$ - and \textbf{Hyperintensionality}, and the component of the argument that recovery is pertinent for simple frameworks is untouched. 

\subsection{Global Consistency Constraints}
By contrast, the problem of rationality is difficult - not because it requires a highly complex formal mechanism to solve, but because there are so many different degrees of rationality one might want to impose on the modelled epistemic agent. 
\par
One main problem of simple frameworks addressed in \S was that simple frameworks allow the belief set to consist both $\phi$ and $\lnot \phi$. However, if a given $\SK$ already consists of both verifier and falsifier of a formula $\phi$ (thus verifiers for $\phi$ and $\lnot \phi$), it would be arbitrary to choose one from $\phi$ or $\lnot \phi$ to include in $\K$. To use the metaphor of ``source'': if an epistemic agent has both sources, or reasons of believing that $\phi$ and believing that $\lnot \phi$, it would seem incorrect to ascribe only one of the belief to the agent. Instead, since we want to model a somewhat ideal epistemic agent, not all sorts of bizarre believer, we could simply forbid the modelled epistemic agent to have both verifier and falsifier for any formula $\phi$ in its ``source''. 
\par
This idea immediately suggests some constraints we can impose on the definition of $\SK$. Since we have two notions of verification and falsification, there are in fact more than one constraint arise from the idea above. I will call the constraints presented in this subsection ``global consistency constraints'', to contrast with the ``local'' consistency constraints which will be presented shortly. 
\begin{definition}[Set of Belief-State Sets with Global Consistency Constraints]
    Given a model $\bm{M}$, 
    \begin{enumerate}
        \item[\emph{(0)}] $\Sigma_0 = \{\SK : \SK \subseteq S \}$
        \item[\emph{(1)}] $\Sigma_1 = \{\SK : \SK \subseteq S \ \& \ \forall s, t \in \SK, \lnot \exists \phi \ (s \ver \phi \ \& \  t \fal \phi) \}$
        \item[\emph{(2a)}] $\Sigma_{2a} = \{\SK : \SK \subseteq S \ \& \ \forall s, t \in \SK, \lnot \exists \phi \ (s \ver \phi \ \& \  t \inefal \phi) \}$
        \item[\emph{(2b)}] $\Sigma_{2b} = \{\SK : \SK \subseteq S \ \& \ \forall s, t \in \SK, \lnot \exists \phi \ (s \inever \phi \ \& \  t \fal \phi) \}$
        \item[\emph{(3)}] $\Sigma_3 = \Sigma_{2a} \cap \Sigma_{2b}$
        \item[\emph{(4)}] $\Sigma_4 = \{\SK : \SK \subseteq S \ \& \ \forall s, t \in \SK, \lnot \exists \phi \ (s \inever \phi \ \& \  t \inefal \phi) \}$
    \end{enumerate}
\end{definition}
It would be easy to check that $\Sigma_4 \subseteq \Sigma_3 \subseteq \Sigma_{2a} \subseteq \Sigma_1 \subseteq \Sigma_0$, and $\Sigma_4 \subseteq \Sigma_3 \subseteq \Sigma_{2b} \subseteq \Sigma_1 \subseteq \Sigma_0$. 
\par
Before we proceed to define our new frameworks utilizing global consistency constraints, we need to first refine our definition of $+$. If we continue to regard $\SK + \sigma = \SK \cup \{\sigma\}$, the result of expansion would not generally be a belief-state set, i.e., member of a $\Sigma_n$, since $\SK$ may contain a falsifier of a formula that $\sigma$ verifies (or vice versa). The modification, however, is rather simple, for we only need to add a condition for an expansion operation to success, as there was a condition for a contraction operation to success in the original AGM framework.
\begin{definition}[Expansion with Consistency Constraints]
    Given a belief-state set $\SK \in \Sigma_n \subseteq \Sigma_0$ and $\sigma \in S$, 
\par
\[
    \SK +_n \sigma = \begin{cases}
        \SK \cup \{\sigma\} & \text{if } \SK \cup \{\sigma\} \in \Sigma_n \\
        \SK & \text{otherwise}
    \end{cases}
    \]
\end{definition}
And contraction is the same with that defined in \textbf{Definition 2.2}. Notice we can use this definition of expansion as the definition of expansion for \textit{all} frameworks, since the expansion will always success with frameworks containing $\Sigma_0$.
\par
With the new conception of expansion with states, we can now further define a suitable notion of revision with states. The idea is that, if one revise with $\sigma$ on $\SK$, we first contract minimal amount of state from $\SK$ to make the expansion with $\sigma$ successful, and then expand with $\sigma$ with the result of the contractions. First, we define conflict relations $\Theta$ that correspond with each consistency constraints: 
\begin{definition}[Conflict Relations Corresponding to Global Consistency Constraints]
    Given a model $\bm{M}$ and $\sigma, \tau$ both in $S$, 
    \begin{enumerate} 
        \item[\emph{(0)}] $\forall \sigma, \tau \ \lnot \Theta_0(\sigma, \tau)$
        \item[\emph{(1)}] $\Theta_1(\sigma, \tau) \iff \exists \phi \ ((\sigma \ver \phi \ \& \  \tau \fal \phi) \lor (\tau \ver \phi \ \& \  \sigma \fal \phi))$
        \item[\emph{(2a)}] $\Theta_{2a}(\sigma, \tau) \iff \exists \phi \ ((\sigma \ver \phi \ \& \  \tau \inefal \phi) \lor (\tau \ver \phi \ \& \  \sigma \inefal \phi))$
        \item[\emph{(2b)}] $\Theta_{2b}(\sigma, \tau) \iff \exists \phi \ ((\sigma \inever \phi \ \& \  \tau \fal \phi) \lor (\tau \inever \phi \ \& \  \sigma \fal \phi))$
        \item[\emph{(3)}] $\Theta_3(\sigma, \tau) \iff \Theta_{2a}(\sigma, \tau) \lor \Theta_{2b}(\sigma, \tau)$
        \item[\emph{(4)}] $\Theta_4(\sigma, \tau) \iff \exists \phi \ ((\sigma \inever \phi \ \& \  \tau \inefal \phi) \lor (\tau \inever \phi \ \& \  \sigma \inefal \phi))$
    \end{enumerate}
\end{definition}
Then we can define revision with states via Quasi-Levi-Identity:
\begin{definition}[Revision with Consistency Constraint]
    Given a model $\bm{M}$ and $\sigma \in S$ \par
    $\SK *_n \sigma = (\SK \setminus \{\tau: \Theta_n (\sigma, \tau)\}) \cup \{\sigma\}.$ 
    \par 
    That is, let $\langle \tau_i \rangle_{i=1}^k$ be an arbitrary enumeration of all states in $\mathcal{S}^\mathcal{K}$ such that $\Theta_n(\sigma, \tau_i)$ holds. We construct a sequence of contractions where $\mathcal{S}_0 = \mathcal{S}^\mathcal{K}$ and $\mathcal{S}_{i+1} = \mathcal{S}_{i} \div \tau_i$. The revision operation is then defined purely in terms of contraction and expansion as:
\begin{equation*}
    \mathcal{S}^\mathcal{K} *_n \sigma = \mathcal{S}_k +_n \sigma
\end{equation*}
\end{definition}
Now we are ready to define new frameworks using global consistency constraint. In principle, we could compose numerous frameworks with the resources given, but only four of them would worth closer examination: 
\begin{definition}[Frameworks $\mathcal{ED}_1, \ \mathcal{ED}_3, \ \mathcal{ED}_4, \text{ and } \mathcal{ID}_4$]
    Given a model $\bm{M}$, \par
     The framework $\mathcal{ED}_1$ is a quintuple $\langle \bm{M}, \Sigma_1, +_1, \div, \Delta_{\mathcal{ED}} \rangle$; \par
    The framework $\mathcal{ED}_3$ is a quintuple $\langle \bm{M}, \Sigma_3, +_3, \div, \Delta_{\mathcal{ED}} \rangle$; \par
    The framework $\mathcal{ED}_4$ is a quintuple $\langle \bm{M}, \Sigma_4, +_4, \div, \Delta_{\mathcal{ED}} \rangle$; \par
    The framework $\mathcal{ID}_4$ is a quintuple $\langle \bm{M}, \Sigma_4, +_4, \div, \Delta_{\mathcal{ID}} \rangle$.
\end{definition}
These (perhaps except $\mathcal{ED}_1$) are the first frameworks presented in the paper which satisfy all the desiderata in \S 1.1. Notice that any $\mathcal{ID}$-series frameworks which has weaker consistency constraint than that of $\Sigma_4$ would still suffer from blatant inconsistency ($\{\phi, \lnot \phi\} \subseteq \K$). Suppose $s \inever p$ and $t \inefal p$, but $p$ is not exactly verified or falsified by any of them. Then under consistency constraint weaker than that of $\Sigma_4$ it could permit $\{s, t\} \subseteq \SK$, and by the definition of $\Delta_{\mathcal{ID}}$, it could therefore permit $\{p, \lnot p\} \subseteq \K$. However, obviously, $\mathcal{ED}_1$ would not suffer from the same problem. 
\par
The difference between three new $\mathcal{ED}$-series frameworks can be appreciated by examining which set of formulas are possible subset of $\K$ under different frameworks. For example, we can observe that if one say $\mathcal{ED}_1$ satisfy \textbf{Local Rationality}, one is saying in a quite weak sense of rationality, since it is possible that $\{p \land \lnot p\} \subseteq \K$, or $\{p \land q, \lnot p \land \lnot q \} \subseteq \K$. Suppose $s_1$ and $s_2$ are sole exact verifier and falsifier of $p$ respectively, and $t_1$ and $t_2$ are sole exact verifier and falsifier of $q$ respectively. Let us further suppose  no mentioned state is a part of another mentioned state, and $p$ and $q$ are only atomic formulas. Then necessary and sufficient condition for $\{p \land q, \lnot p \land \lnot q \} \subseteq \K$ to be the case under $\mathcal{ED}_1$ is that $\{s_1 \sqcup t_1, \ s_2 \sqcup t_2 \} \subseteq \SK$. It should be noticed that there are no formulas \textit{exactly} verified by $s_1 \sqcup t_1$ yet \textit{exactly} falsified by $s_2 \sqcup t_2$, and also vice versa: atomic formulas $p$ or $q$ (and therefore $\lnot p$ or $\lnot q$) are merely inexactly verified/falsified by them; $p \lor q$ (and therefore $\lnot (p \lor q) \equiv (\lnot p \land \lnot q)$) is exactly falsified (verified) by $s_2 \sqcup t_2$ yet merely inexactly verified (falsified) by $s_1 \sqcup t_1$; and $p \land q$ (and therefore $\lnot (p \land q) \equiv (\lnot p \lor \lnot q)$) is exactly verified (falsified) by $s_1 \sqcup t_1$ yet merely inexactly falsified (verified) by $s_2 \sqcup t_2$. 
\par
$\{p \land q, \lnot p \land \lnot q \}$, however, is not a possible subset of $\K$ under $\mathcal{ED}_3$, and \textit{a fortiori} $\mathcal{ED}_4$. For any two states $\sigma$ and $\tau$, if $\sigma \ver p \land q$, $\sigma$ and $\tau \ver \lnot p \land \lnot q$, $\tau$ would have an exact verifier of $\lnot p$ as part (possibly improper), which is an exact falsifier of $p$, and therefore an exact falsifier of $p \land q$. Thus $\tau$ would have an exact falsifier of $p \land q$ as part, i.e., is an inexact falsifier of $p \land q$. Therefore, any such set $\{\sigma, \tau\} \nsubseteq \SK$ under $\mathcal{ED}_3$ and  $\mathcal{ED}_4$, and therefore $\{p \land q, \lnot p \land \lnot q \} \subseteq \K$ is also impossible under those frameworks. 
\par
Likewise, we might wish find a set of formulas that is possibly a subset of $\K$ under $\mathcal{E}_3$ but not under $\mathcal{E}_4$. This is however, not as simple as the previous case. For, there is no purely syntactical way to define a possible $\K$-subset under $\mathcal{E}_3$ which is not a possible $\K$-subset under $\mathcal{E}_4$ (proof in appendix). But still we can spot a difference between those two frameworks somewhat simillarly, by getting a little help with the semantics. 
\par



% --- BIBLIOGRAPHY ---
\bibliography{refs} % Make sure 'refs.bib' exists in the folder
\end{document}